

\textbf{Определение.} 
\textit{Количество способов выбрать $k$ предметов из $n$ различных предметов обозначается $C_n^k$ (читается «це из эн по ка»)}

\textit{Найдем формулу для $C_n^k$. Первый предмет можно выбрать $n$ способами, второй $n-1$ способом и так далее, последний $k$-й элемент можно выбрать уже ( $n-k+1$ ) способом. Однако один и тот же набор из $k$ предметов мы тут выбирали по-разному - в разном порядке. Всего способов упорядочить $k$ предметов - $k!=k *(k-1) *(k-2) * \ldots * 1$. Получаем формулу}

$$
C_n^k=\frac{n \cdot(n-1) \cdot(n-2) \ldots \cdot(n-k+1)}{k \cdot(k-1) \cdot \ldots \cdot 2 \cdot 1}
$$


Правую часть можно преобразовать, умножив и числитель, и знаменатель на $(n-k)!=(n-k) *(n-k+1) * \ldots * 1$.

$$
\frac{n \cdot(n-1) \cdot \ldots \cdot(n-k+1)}{k \cdot(k-1) \cdot \ldots \cdot 2 \cdot 1}=\frac{n \cdot(n-1) \cdot \ldots \cdot(n-k+1) \cdot(n-k) \cdot(n-k-1) \ldots \cdot 2 \cdot 1}{k \cdot(k-1) \cdot \ldots \cdot 2 \cdot 1 \cdot(n-k) \cdot(n-k-1) \ldots \cdot 2 \cdot 1}=\frac{n!}{k!(n-k)!}
$$


Так мы получаем две формулы для $C_n^k$, полезно запомнить обе - иногда удобнее одна, иногда другая.

$$
C_n^k=\frac{n \cdot(n-1) \cdot(n-2) \ldots \cdot(n-k+1)}{k \cdot(k-1) \cdot \ldots \cdot 2 \cdot 1}=\frac{n!}{k!(n-k)!}
$$


\problem{1}
Сколькими способами можно выбрать 4 конфеты из имеющихся 7 различных?

\textbf{Ответ:} $C_7^4=35$.

\textbf{Решение.} Нам нужно выбрать 4 предмета из 7, это можно сделать $C_7^4$ способами.

\problem{2}
На окружности отмечено 10 точек. Сколько существует треугольников с вершинами в этих точках?

\textbf{Ответ:} $C_{10}^3=120$ треугольников.

\textbf{Решение.} Треугольник задается своими вершинами. Т.е. нам нужно выбрать 3 точки из 10, это можно сделать $C_{10}^3$ способами.

\problem{3}
У семиклассника Лени есть 7 детективов, а у восьмиклассника Бори 8 книг по математике. Сколькими способами они могут обменять три книги одного на три книги другого?

\textbf{Ответ:} $C_7^3 \cdot C_8^3=1960$ способа.

\textbf{Решение.} Каждую из $C_7^3$ троек книг Лени можно обменять на любую из $C_8^3$ троек книг Бори, т.е. всего $C_7^3 \cdot C_8^3$ способов обмена.

\problem{4}
Из двух математиков и десяти экономистов надо составить комиссию из семи человек. Сколькими способами можно составить комиссию, если в нее должен входить хотя бы один математик?

\textbf{Ответ:} $C_{10}^5+2 \cdot C_{10}^6=672$.

\textbf{Решение 1.}
Разобьем все способы на две непересекающиеся группы: 1) в комиссии один математик 2) в комиссии два математика. И посчитаем количество способов в каждой группе отдельно. Если оба математика вошли в комиссию, то надо еще подобрать 5 экономистов. Так как всего экономистов 10 - это можно сделать $C_{10}^5$ способами. Если же в комиссию входит только один математик, то нужно еще подобрать 6 экономистов - это можно сделать $C_{10}^6$ способами. Так как выбрать одного математика можно 2 способами, то всего комиссий с одним математиком возможно $2 \cdot C_{10}^6$.

\textbf{Решение 2.}
Посчитаем способы выбора 7 человек для комиссии независимо от их профессий, затем вычтем те способы в которых нет математиков. Всего 12 человек (2 математика и 10 экономистов) выбрать из них 7 человек можно $C_{12}^7$ способами. Набрать комиссию из одних экономистов можно $C_{10}^7$ способами. Всего способов набрать комиссию с математиками $C_{12}^7-C_{10}^7= 792-120=672$.

\problem{5}
Из 12 эльфов и 10 гномов выбирают команду, состоящую из пяти человек для борьбы со Злом. Сколькими способами можно выбрать эту команду так, чтобы в нее вошло не более трех гномов?

\textbf{Ответ:}
$C_{10}^3 \cdot C_{12}^2+C_{10}^2 \cdot C_{12}^3+C_{10}^1 \cdot C_{12}^4+C_{12}^5=23562$.

\textbf{Решение 1.}
Аналогично решению $1 к$ предыдущей задаче разобьем все способы на группы: в команду входит 3 гнома, в команду входит 2 гнома, в команду входит 1 гном, в команду не входит ни одного гнома.
Посчитаем количество способов в каждом из случаев.
Если гномов три, то этих трех гномов можно выбрать $C_{10}^3$ способами. Оставшихся двух эльфов можно подобрать $C_{12}^2$ способами. Значит, всего в этом случае будет $C_{10}^3 \cdot C_{12}^2$ способов.
Аналогично, способов, когда гномов в команде два будет $C_{10}^2 \cdot C_{12}^3$. Способов, когда в команде один гном будет $C_{10}^1 \cdot C_{12}^4$, и, наконец, когда гномов в команде нет, будет $C_{12}^5$ способов. Итого всего $C_{10}^3 \cdot C_{12}^2+C_{10}^2 \cdot C_{12}^3+C_{10}^1 \cdot C_{12}^4+C_{12}^5=23562$ способов.

\textbf{Решение 2.}
Аналогично решению 2 к предыдущей задаче сосчитаем способы набора команды не учитывая расовую принадлежность и выкинем способы в которых 4 и 5 гномов. Всего претендентов 22 (12 эльфов и 10 гномов), значит всего способов набрать команду $C_{22}^5$. Вариантов команд в которых все гномы $C_{10}^5$, а вариантов команд в которых ровно 4 гнома $C_{10}^4 \cdot C_{12}^1$. Всего нужных нам вариантов команд: $C_{22}^5-C_{10}^5-C_{10}^4 \cdot C_{12}^1=26334-252-2520=23562$.

\problem{6}
Как известно, для участия в лотерее «Спортлото» нужно указать шесть номеров из имеющихся на карточке 45 номеров.

a) Сколькими способами можно заполнить карточку «Спортлото»?

б) После тиража организаторы лотереи решили подсчитать, каково число возможных вариантов заполнения карточки, при которых могло быть угадано ровно три номера. Помогите им в этом подсчете.


\textbf{Ответ:} а) $C_{45}^6=8145060$; б) $C_6^3 \cdot C_{39}^3=182780$.

\textbf{Решение.} а) нам нужно отметить какие-то 6 из 45 номеров. Это можно сделать $C_{45}^6$ способами.

б) Нам даны выигравшие 6 номеров. Нужно узнать количество способов заполнить карточку, так, чтобы было угадано ровно три из них. Чтобы получить такую карточку, нужно выбрать 3 выигрышных номера (из 6) и 3 невыигрышных (из 39). Это можно сделать $C_6^3 \cdot C_{39}^3$ способами.

\problem{7}
На дежурстве в столовой два школьника должны разносить стаканы, трое разливать компот, пятеро - разносить второе. Сколькими способами десять школьников могут разделиться на дежурстве?

\textbf{Ответ:} $C_{10}^2 \cdot C_8^3=2520$.

\textbf{Решение.} Двоих школьников на разнос стаканов можно выбрать $C_{10}^2$ способами. Для каждого из этих способов выбрать из оставшихся школьников еще троих для разливания компота можно $C_8^3$ способами. Т.е. всего способов выбрать двоих для стаканов и троих для компота $C_{10}^2 \cdot C_8^3$. Выбрать из оставшихся пятерых пять для второго можно одним способом, т.е. всего $C_{10}^2 \cdot C_8^3=\frac{10!}{2!8!} \cdot \frac{8!}{3!5!}=\frac{10!}{2!3!5!}=2520 \quad$ способов разделиться на дежурство десяти школьникам.
Замечание. Обязанности можно было бы распределять в другом порядке. Например, сначала направить 3 людей на компот - это можно сделать $C_{10}^3$. Потом отправить 5 человек разносить второе - это можно сделать $C_7^5$ способами. А оставшихся двоих отправить расставлять стаканы. Значит всего способов распределить детей $C_{10}^3 \cdot C_7^5=\frac{10!}{3!7!} \cdot \frac{7!}{5!2!}=\frac{10!}{2!3!5!}$. Как видно, после перемножения ответ получился такой же, как и раньше.

Еще задача.

\problem{8}
Найдите среднее арифметическое всех девятизначных чисел составленных из 3 троек и 6 шестерок.

\textbf{Ответ:} 555555555 (число, состоящее из 9 пятерок).

\textbf{Решение.} Среднее арифметическое это сумма всех чисел поделить на их количество. Чтобы узнать количество чисел необходимо сосчитать сколькими способами можно из трех троек и шести шестерок составить девятизначное число. Для этого достаточно из 9 мест числа выбрать, например, 3 для троек. Это можно сделать $C_9^3=84$ способами, значит и чисел будет 84. Теперь найдем сумму этих 84 чисел, сделаем это поразрядно. Найдем сумму всех последних цифр наших 84 чисел. В конце $y$ числа может стоять тройка - ровно у $C_8^2=28$ чисел, а шестерка - y $C_8^3=56$ чисел. Значит, сумма последних разрядов всех чисел равна $3 * 28+6 * 56=420$. Посчитаем теперь сумму предпоследних разрядов. Опять же там будет 28 троек и 56 единиц, но так как это разряд десяток то вклад в сумму будет равен $(3 * 28+6 * 56) * 10=420 * 10$. Аналогично вклад сотен равен $420 * 100 u$ так далее. Значит, сумма всех чисел равна $420 *(1+10+100+\ldots+100000000)=420 * 111111111$. Чтобы найти среднее арифметическое нужно найденную сумму разделить на 84. Среднее арифметическое равно: $420 * 111111111 / 84=55555555$ (число из 9 пятерок).